\section{Training and Evaluation Protocol}
\label{sec:protocol}

\subsection{Training}
\label{sec:training}

The student is Qwen2.5-1.5B-Instruct. We use the plain Hugging Face \code{Trainer} with the loss masked by the token length of
the prompt, rather than TRL's \code{SFTTrainer}, whose completion-only collator searches for the assistant marker (which
appears twice in answer-phrasing samples) and silently masks a whole sample when it does not match. Hyperparameters follow
Octopus~v2: AdamW, learning rate $5\times10^{-5}$, 10 warm-up steps, linear decay, 3 epochs, equal loss weights for the new
tokens (Octopus~v2 found weighting unnecessary for few tokens), effective batch 16, maximum length 256, bf16. LoRA uses rank 16,
$\alpha=32$ on \code{q,k,v,o,up,down} projections; because the new tokens must be learned, \code{embed\_tokens} and
\code{lm\_head} are trained in full. Qwen2.5-1.5B ties these two matrices; saving a merged PEFT model with tied weights
silently discards the trained output layer, so we untie them before training. With the two full matrices, 480M of 2.26B
parameters are trainable (21\%). Each configuration is trained with seeds 42 and 7 on a rented A100 40\,GB, L40S or L4
(v6.1: about 57 minutes and 6,288 steps per seed). A checkpoint is saved per epoch and chosen on 600 stratified validation
samples: the highest accuracy on clear commands among checkpoints with false refusals $\le5\%$, earlier epoch on ties.

\subsection{Scoring}
\label{sec:scoring}

All models, including prompted baselines whose JSON output is converted to the functional format, are scored by the same
code. An output is \emph{correct} if it matches the reference after normalization (``rule C''): arguments equal to their
schema default are dropped on both sides (\code{get\_odometer()} = \code{get\_odometer(unit="km")}); for ambiguous items only
the function and the direction are checked (on/off, open/close, up/down, sign of a delta), and a proposal may pick any one of
several acceptable functions; suggesting the high beam is always wrong. Real behavioral differences (delta 1 vs.\ 2, miles vs.\
kilometres) are never forgiven. Rule C was fixed before round R2 and old results were re-scored with it; under the earlier
exact-match rule the same R1 outputs score 4.1 points lower on clear commands.

Errors are classified into seven types: wrong function, wrong argument, false refusal, over-triggering (calling a real
function on an out-of-scope request), unnecessary clarification, missing clarification, and unparseable output. We call
over-triggering, wrong function, wrong argument and missing clarification \emph{wrong actions}: the car would do something the
user did not ask for. We report per-group accuracy and never a single pooled number: a model that refuses everything already
scores 51\% on the clean test.

\subsection{Statistical protocol}
\label{sec:stats}

The test sets are small (74 clear commands; one item is 1.35 points), so we use Wilson intervals \citep{wilson1927} and the
exact two-sided McNemar test \citep{mcnemar1947} on paired outputs. McNemar accounts for the choice of test items but not for
training randomness: when we retrained round R2 with a second seed on identical data, 33 of 230 items changed outcome and
McNemar between the two seeds gave $p=0.005$. From then on every configuration was trained with two seeds (three for the final one, Appendix~\ref{app:seeds}), results are the
mean of the two, and a difference is claimed only if the mean difference exceeds the spread between seeds and, for paired
tests, if it is significant for both seeds in the same direction (one seed: ``tie or slightly better''). The criteria for
accepting each round, choosing thresholds and compression levels were recorded before the corresponding test run, and we
did not change them after seeing scores. Two seeds remain few; our conclusions are correspondingly modest.
